\section{Marginal axis effects}
\label{sec:marginals}

This section reports the pooled outcome means for every level of every axis. Marginal
means average completed seeds within a design cell and then average design cells within an
axis level, so that better-replicated configurations do not receive extra weight.
Section~\ref{sec:contrasts} repeats the same axes as matched within-design contrasts, which
is the stronger read; the marginals are reported first because they establish the scale of
every quantity.

Table~\ref{tab:headline-axes} gives the pooled means over all four benchmarks.
Appendix~\ref{app:marginals} gives the benchmark-resolved version with bootstrap intervals
for all thirteen metrics.

\intable{tab_headline_axes}

\subsection{Capacity}

Accuracy rises monotonically with parameter count: 0.356 at 0.6B, 0.497 at 1.7B, 0.651 at
4B, 0.672 at 8B, 0.723 at 14B and 0.727 at 32B. The gain from 0.6B to 14B is 0.367
accuracy points, larger than the effect of any other axis in the sweep by a factor of
roughly four. The gain from 14B to 32B is 0.004, which is within noise and is the plateau
discussed in Section~\ref{sec:scaling}.

The capacity effect on calibration is signal-dependent and this is the first place the
three-signal decomposition earns its keep. Vote-fraction calibration error improves
sharply with scale on every benchmark. Comparing 32B against 0.6B in matched contrasts,
vote calibration error falls by 0.233 on GPQA, 0.321 on MMLU-Pro, 0.387 on TruthfulQA and
0.068 on MATH. Per-agent calibration error does not follow: over the same contrast it is
essentially unchanged on GPQA at $+0.004$, worsens substantially on MMLU-Pro at $+0.234$,
and improves on TruthfulQA at $-0.148$.

The larger model therefore makes agreement a more honest signal much more reliably than it
makes an individual model's own probability more honest. Section~\ref{sec:calibration}
takes this further and shows that the two move at different rates with scale, which
reverses the direction the sweep was designed to test.

\subsection{Reasoning budget}

Accuracy rises with the reasoning rung: 0.522 with thinking disabled, 0.569 at b512, 0.620
at b2048, 0.633 at b8192 and 0.634 at unlimited. The pattern is a large first step, a
second substantial step, then saturation. The final rung adds nothing measurable over
b8192 in the pooled mean, and Section~\ref{sec:scaling} shows it is zero or negative on
every benchmark in matched contrasts.

The calibration pattern splits by signal in the opposite direction to capacity. Moving
from disabled to unlimited thinking, vote calibration error improves on all three
multiple-choice benchmarks, by 0.058 on GPQA, 0.102 on MMLU-Pro and 0.049 on TruthfulQA,
and the vote delta falls by 0.098, 0.211 and 0.074 respectively. Per-agent calibration
error moves the other way on all three: $+0.041$ on GPQA, $+0.110$ on MMLU-Pro and
$+0.025$ on TruthfulQA.

Reasoning also reduces signed overconfidence. The dose-response slopes per unit of
$\log(1 + \text{realized reasoning tokens})$ are $-0.0082$ on GPQA, $-0.0140$ on MMLU-Pro
and $-0.0058$ on TruthfulQA for the signed confidence gap, against accuracy slopes of
$+0.0106$, $+0.0175$, $+0.0153$ and $+0.0106$ on the four benchmarks. Per-agent
calibration error slopes are simultaneously positive at $+0.0045$, $+0.0126$ and
$+0.0034$.

The combination is coherent and is one of the more interesting results in the sweep: extra
thinking makes models more accurate and less overconfident on average, while making their
raw option-logprob distribution a worse-calibrated probability. The natural reading is that
reasoning shifts the confidence distribution past the calibration point rather than onto
it, replacing overconfidence with underconfidence. Section~\ref{sec:calibration} examines
the reliability diagrams and the Cox recalibration slopes that bear on this.

Figure~\ref{fig:token-dose} shows the dose relationship directly.

\begin{figure}[H]
\centering
\includegraphics[width=\linewidth,height=0.28\textheight,keepaspectratio]{fig_token_dose.pdf}
\caption{Accuracy against realized reasoning tokens, by capacity and benchmark. The
horizontal axis is the number of thinking tokens the model actually spent, not the budget
it was given.}
\label{fig:token-dose}
\end{figure}

The distinction between the nominal rung and the realized dose matters at 32B, where the
context cap prevents the model from spending its allowance. At b8192 the median 32B cell
spends 3{,}472 thinking tokens against 3{,}981 for 14B; at unlimited it spends 3{,}124
against 4{,}731. The largest model is the one least able to use the largest budget.

\subsection{Prompt complexity}

Accuracy is 0.573 at prompt level 0 and 0.617 at level 3, a pooled gain of 0.044. The
matched contrast is positive on every benchmark: $+0.015$ on GPQA, $+0.060$ on MMLU-Pro,
$+0.070$ on TruthfulQA and $+0.018$ on MATH.

Prompt scaffolding also improves agreement calibration everywhere: vote calibration error
falls by 0.032, 0.050, 0.059 and 0.008 on the four benchmarks, and the vote delta falls by
0.021, 0.083 and 0.077 on the three where it is defined. Per-agent calibration error is
mixed, improving by 0.011 on GPQA but worsening by 0.033 on MMLU-Pro and 0.018 on
TruthfulQA.

This is the most uniformly favourable axis in the sweep on the two headline outcomes, and
it is also the cheapest, since the L3 template adds a few hundred prompt tokens and no
additional generation. The caveat is the same one that appears on the reasoning axis: the
intervention that makes the system better at answering makes the individual model's raw
confidence less calibrated, even though the L3 template explicitly instructs the model to
calibrate.

\subsection{Context sharing}

Pooled accuracy is 0.585 for artifact-only and 0.614 for plus-CoT. The axis is identified
only within the two communicating topologies, since the other two have no peer context. In
matched contrasts the accuracy gain is $+0.016$ on GPQA, $+0.016$ on MMLU-Pro, $+0.009$ on
TruthfulQA and $+0.052$ on MATH.

The calibration cost is consistent in sign across all four: vote calibration error rises by
0.037, 0.030, 0.018 and 0.005, and the vote delta rises by 0.036, 0.023 and 0.016 on the
three benchmarks where it is defined. Per-agent calibration error barely moves at
$+0.000$, $+0.007$ and $+0.003$.

Exposing peer rationales therefore helps the system answer and harms the trustworthiness
of its agreement, without measurably changing how well calibrated the individual agents
are. This is precisely the signature of induced error correlation: the agents are
individually no worse calibrated, but they agree more often for reasons that are not
evidence.

Section~\ref{sec:contrasts} shows that the accuracy benefit, unlike the calibration cost,
is regime-dependent and largely disappears in the large-capacity high-reasoning corner.

\subsection{Topology}

Table~\ref{tab:topology-means} gives the complete outcome profile of the four topologies.

\intable{tab_topology_means}

Pooled accuracy is 0.564 for single agent, 0.598 for independent, 0.615 for decentralized
and 0.587 for centralized. Decentralized debate has the highest average accuracy and the
lowest error amplification at $A_e = 0.886$; centralized orchestration is the weakest of
the three multi-agent conditions on average accuracy despite consuming the same six
generation turns as decentralized.

Coordination efficiency ranks in the opposite order to accuracy, because it is accuracy
divided by relative turn count: 0.564 for single agent by construction, 0.198 for
independent, 0.102 for decentralized and 0.097 for centralized. Turn overhead is 200
percent for independent and 500 percent for both two-round topologies. Message counts are
zero for the two non-communicating conditions and six for both communicating conditions.
The cost proxy runs from 15{,}339 for single agent to 56{,}177 for independent, 102{,}771
for centralized and 120{,}897 for decentralized.

Every multi-agent condition has $A_e < 1$, meaning all three make fewer errors than a
matched single agent. The ordering, 0.886 for decentralized, 0.928 for centralized and
0.934 for independent, says debate suppresses the largest share of single-agent errors.
Section~\ref{sec:items} shows this coexists with debate producing more unanimous incorrect
consensus, which is not a contradiction: the two statistics count different events.

\subsection{Agent count}

The one-against-three contrast averages the available multi-agent topologies against the
matched single-agent baseline at artifact-only context. Accuracy gains are $+0.014$ on
GPQA, $+0.031$ on MMLU-Pro, $+0.019$ on TruthfulQA and $+0.056$ on MATH, all with
intervals excluding zero. The cost proxy rises by 84{,}758, 58{,}360, 35{,}356 and 62{,}219
respectively, and error amplification falls by 0.025, 0.093, 0.066 and 0.137.

The accuracy gains are real and consistently signed, but they are small relative to the
cost multiplier of three to six, and they are three to four times larger on the
mathematical benchmark than on the two knowledge benchmarks. As stated in
Section~\ref{sec:design}, this contrast compares system families rather than agent counts:
the three-agent conditions differ from the single-agent condition in agent count,
coordination structure, round count and the presence of self-consistency sampling
simultaneously.

\subsection{Complete marginal tables}

Appendix~\ref{app:marginals} contains the benchmark-resolved marginal means with
hierarchical bootstrap intervals for accuracy, per-agent calibration error, vote
calibration error, final-producer calibration error, both deltas, both signed gaps, the
cost proxy, mean total tokens, mean reasoning tokens, coordination efficiency and error
amplification. That is thirteen tables covering 987 estimates, which together constitute
the complete marginal description of the sweep.
